\section{Preliminaries}\label{section2}
In this section we review some of the models for the space of orthogonal complex structures on a Euclidean space following~\cite[Section~IV.9]{LM}, to which we refer for further details.

Let $(V, (\cdot, \cdot ))$ be a $2n$-dimensional real inner product space. We will write \[J(V)=\big\{ J \in \End(V) \colon J^2=-I, \, J \text{ is orthogonal} \big\}\]
for the space of orthogonal complex structures on $V$. This space naturally has the structure of a complex algebraic variety, given by the isomorphism
\[ J(V) \xrightarrow{\phi} \Gr_n^{\rm Iso}(V\otimes \C) \]
with the Grassmannian of $n$-dimensional complex subspaces of $V\otimes \C$, which are isotropic with respect to the complex bilinear form on $V\otimes \C$ obtained from the inner product on $V$ by linear extension. The map $\phi$ is defined by
\[ \phi(J) = \{ v\otimes 1 + Jv \otimes {\rm i} \colon v \in V\} \]
(it assigns to $J$ the plane $V_J^{0,1} \subset V\otimes \C$, which is the $-{\rm i}$ eigenspace of the complex linear extension of $J$ to $V\otimes \C$).

There is a tautologous complex vector bundle on $J(V)$ defined by
\[J(V) \times V \xrightarrow{\pi_1} J(V),\]
where the fiber over $J\in J(V)$ is given the complex structure determined by $J$.
The isomorphism $\phi$ is covered by a bundle map to the dual (or conjugate) tautological bundle over the Grassmannian.

\begin{Lemma}\label{positchern}
Consider the universal bundle
\[
J\big(\R^{2n}\big)={\rm O}(2n)/{\rm U}(n) \xrightarrow{\iota} {\mathcal B}{\rm U}(n) \to {\mathcal B}{\rm O}(2n).
\]
Then
\begin{enumerate}\itemsep=0pt
\item[$(i)$] $\iota$ classifies the tautologous complex vector bundle on $J\big(\R^{2n}\big)$.
\item[$(ii)$] $\iota^*(c_1)$ generates a subgroup of index $2$ in the infinite cyclic second cohomology group of
each of the two components of $J\big(\R^{2n}\big)$.
\item[$(iii)$] $\iota^*(c_1)$ is positive. In particular, $\iota^*(c_1)^{\frac{n^2-n}{2}}$ evaluates on the orientation class of each of the two components of $J\big(\R^{2n}\big)$ to a positive integer.
\end{enumerate}
\end{Lemma}
\begin{proof}
$(i)$ The universal bundle can be realized by taking ${\mathcal B}{\rm U}(n)={\mathcal E}{\rm O}(2n)/{\rm U}(n)$. The universal bundle over ${\mathcal B}{\rm U}(n)$ is then given by ${\mathcal E}{\rm O}(2n)\times_{{\rm U}(n)} \R^{2n} \to {\mathcal E}{\rm O}(2n)/{\rm U}(n)$, and pulls back along the inclusion of the fiber $\iota$ to ${\rm O}(2n) \times_{{\rm U}(n)} \R^{2n} \to {\rm O}(2n)/{\rm U}(n)$, which is the homogeneous expression of the tautologous bundle.

$(ii)$ The inclusion-induced map ${\mathcal B}{\rm U}(n) \to {\mathcal B}{\rm O}(2n)$ factors through the
double cover ${\mathcal B}{\rm SO}(2n) \allowbreak \to {\mathcal B}{\rm O}(2n)$, and the two components of ${\rm O}(2n)/{\rm U}(n)$ are,
respectively, the fiber ${\rm SO}(2n)/{\rm U}(n)$ of ${\mathcal B}{\rm U}(n) \to {\mathcal B}{\rm SO}(2n)$ and the image
of this fiber under a lift to ${\mathcal B}{\rm U}(n)$ of the nontrivial deck transformation of
${\mathcal B}{\rm SO}(2n)$. Therefore it suffices to prove the statement for the component
${\rm SO}(2n)/{\rm U}(n)$. This follows from the Serre spectral sequence of the bundle
${\mathcal B}{\rm U}(n) \to {\mathcal B}{\rm SO}(2n)$ (see for instance~\cite[Theorem~III.6.11]{MiTo}).

$(iii)$ $\iota^*(c_1)$ is the pullback of $c_1$ of the dual tautological bundle under the embedding
\[J\big(\R^{2n}\big) \xrightarrow{\cong} \Gr_n^{\rm Iso}\big(\C^{2n}\big) \subset \Gr_n\big(\C^{2n}\big).\]
The $n^{\mathrm{th}}$ exterior power of the dual tautological bundle on $\Gr\big(\C^{2n}\big)$ gives the
Pl\"ucker embedding of the Grassmannian in projective space. Hence $\iota^*(c_1)$ is the first
Chern class of a line bundle on $J\big(\R^{2n}\big)$ whose sections embed $J\big(\R^{2n}\big)$ in projective space and is therefore positive.\qedhere
\end{proof}

The above identification of orthogonal complex structures with isotropic planes in
$V\otimes \C$ leads to another description of the space $J(V)$ which will play an
important role in this paper.


Let $\Cl(V)$ denote the Clifford algebra determined by the quadratic form $q(v)=-\|v\|^2$, and let $\CCl(V)=\Cl(V)\otimes_\R \C$
denote its complexification. Let $S_{\C}$ denote an irreducible ($2^n$-dimensional) $\CCl(V)$-module.

For each element, i.e., \emph{spinor}, $\sigma \in S_{\C}$ we have the map $V\otimes_\R \C \to S_{\C}$ given by right Clifford multiplication by $\sigma$. A spinor is said to be \emph{pure} if the associated kernel is half-dimensional. Let ${\mathcal P}S _{\C} \subset S_\C$ denote the subset of pure spinors. As the subspace $\ker (\cdot \sigma) \subset V\otimes \C$ is isotropic
there is a natural map from the projectivization of the set of pure spinors to the isotropic Grassmannian
\[
\mathbb{P}( {\mathcal P}S _\C ) \xrightarrow{\psi} \Gr^{\rm Iso}_{2n}(V\otimes \C)
\]
defined by $[\sigma] \mapsto \ker (\cdot \sigma)$. This is a map of algebraic varieties and is in fact an ${\rm SO}(2n)$-equivariant isomorphism~\cite[Proposition~IV.9.7]{LM}.

The space $J(V)$ has two connected components corresponding to the two possible orientations induced by the complex structure. A choice of orientation for $V$ leads to a decomposition
\[J(V) = J_+(V) \coprod J_-(V)\]
with $J_+(V)$ the component corresponding to the given orientation. In turn this decomposes the Grassmannian of isotropic subspaces into the components of positive and negative isotropic subspaces.

On the level of spinors this decomposition takes the following form.
The complex volume element in $\CCl(V)$ is the element
$\omega_\C = {\rm i}^n e_1 \cdots e_{2n}$, where
$\{e_i\}$ is any oriented orthonormal basis. As $\omega_{\C}^2=1$, the volume element decomposes
the module $S_C$ into a direct sum of $+1$ and $-1$ eigenspaces
\[
S_\C = S_\C^+ \oplus S_\C^-.
\]
The elements of the summands are called the positive and negative spinors, respectively.
Note that, since $\omega_\C$ anti-commutes with the action of $V \subset \CCl(V)$, both
$S_\C^+$ and $S_\C^-$ are ${\rm spin}^{\rm c}$ representations, called the positive and negative spinor representations, respectively.

Every pure spinor is necessarily either positive or negative and the isomorphisms $\phi$, $\psi$ give rise to isomorphisms of algebraic varieties
\[
\mathbb{P}\big({\mathcal P}S _\C^\pm\big) \cong J_\pm(V).
\]
In dimensions $2n=4,6$ all nonzero spinors are pure~\cite[Remark~IV.9.12]{LM} and the isomorphisms identify the spaces
of orthogonal complex structure with the disjoint union of the projectivization of positive
and negative spinors.

A point $J \in J(V)$ gives rise to a useful description (cf.~\cite[Chapter I, equation~(5.27)]{LM})
of the irreducible module $S_\C$.
We will write $\langle \cdot, \cdot \rangle$ for the Hermitian inner product on the complex vector space $(V,J)$ whose real part is $ (\cdot, \cdot )$.
Let $\Lambda_\C^*(V)$ denote the exterior algebra of the complex vector space $(V,J)$.
Contraction of a vector $v \in V$ with $\omega \in \Lambda_\C^k(V)$ is determined by the expression
\[
v \lrcorner (w_1 \wedge \cdots \wedge w_k) = \sum_{j=1}^k (-1)^{j+1}\langle v, w_j \rangle w_1 \wedge \cdots \wedge \widehat{w_j} \wedge \cdots \wedge w_k.
\]
The action $V \otimes_\R \Lambda_\C^*(V) \to \Lambda_\C^*(V)$ defined by
\begin{equation*}
v \cdot \omega = v \wedge \omega - v \lrcorner \, \omega
\end{equation*}
satisfies
\begin{equation*}
v\cdot(v \cdot \omega) = - \|v\|^2 \omega,
\end{equation*}
as one easily checks using an orthonormal basis for $V$ having a
real multiple of a nonzero $v$ as its first element. The complex linear extension of this action gives $\Lambda_\C^*(V)$ the structure of a $\CCl(V)$-module of complex dimension $2^n$. This is the dimension of the irreducible complex $\CCl(V)$-module, so $\Lambda_\C^*(V)$ is a convenient description of this module.

\iffalse
Indeed, one can consider an orthonormal basis $v_1,Jv_1, \ldots, v_n, Jv_n$ for $V$
with $v=\alpha v_1$ and check that \eqref{clifalg} holds for each $\omega$ in the standard orthonormal basis for $\Lambda_\C^*(V)$:
$$
v \cdot v_{i_1} \wedge \cdots \wedge v_{i_k} = \begin{cases}
- \alpha \ v_{i_2} \wedge \cdots \wedge v_{i_k} & \text{ if } i_1=1\\
\alpha \ v_{1} \wedge v_{i_1} \wedge \cdots \wedge v_{i_k} & \text{ if } i_1 \neq 1
\end{cases}
$$
and hence
$$
v \cdot( v \cdot v_{i_1} \wedge \cdots \wedge v_{i_k} ) =
- \alpha^2 \ v_{i_1} \wedge v_{i_2} \wedge \cdots \wedge v_{i_k} = - \|v\|^2 v_{i_1} \wedge \cdots \wedge v_{i_k}
$$
\fi

The following is an easy computation which we include for convenience.

\begin{Lemma}\label{plusmin}
Let $(V,J)$ be an oriented $2n$-dimensional Euclidean real vector space with an orthogonal complex structure $J$ compatible with the orientation and let $S_\C=\Lambda^*_\C(V)$ be the irreducible $\CCl(V)$-module described above. Then
\[
S_\C^+ = \bigoplus_{k \text{ even }} \Lambda_\C^k(V), \qquad S_\C^- = \bigoplus_{k \text{ odd }} \Lambda_\C^k(V).
\]
\end{Lemma}
\begin{proof}
Let $e_1, \ldots, e_n$ be an orthonormal basis for $(V,J)$. Then we can write the complex volume element as
\[
\omega_\C = {\rm i}^n e_1 Je_1 \cdots e_n Je_n .
\]
Consider the basis $e_{i_1} \wedge \cdots \wedge e_{i_k}$ for $\Lambda^k_{\C}$ with $1\leq i_1 < \cdots < i_k \leq n$. Let $1 \leq m \leq n$
and assume that $m \not \in \{i_1, \ldots, i_k \}$. Then
\[
\begin{aligned}
(e_m Je_m) \cdot e_{i_1} \wedge \cdots \wedge e_{i_k} & = e_m \cdot \big( Je_m \wedge e_{i_1} \wedge \cdots \wedge e_{i_k} \big) \\
& = -\big\langle e_m, Je_m \big\rangle e_{i_1} \wedge \cdots \wedge e_{i_k} = -i \ e_{i_1} \wedge \cdots \wedge e_{i_k}.
\end{aligned}
\]
On the other hand, if $m \in \{i_1, \ldots, i_k\}$, let $l$ be such that $m=i_l$. Then
\[
\begin{aligned}
(e_m Je_m) \cdot e_{i_1} \wedge \cdots \wedge e_{i_k} & = e_m \cdot \big( {-} (-1)^{l-1} \big\langle Je_m, e_{i_l} \big\rangle e_{i_1} \wedge \cdots \wedge \widehat{e_{i_l}} \wedge \cdots \wedge e_{i_k} \big)\\
& = e_m \wedge\big( (-1)^l (-{\rm i}) e_{i_1} \wedge \cdots \wedge \widehat{e_{i_l}} \wedge \cdots \wedge e_{i_k} \big) \\
& = (-1)^{l-1} (-1)^l(-{\rm i}) e_{i_1} \wedge \cdots \wedge e_{i_k} = {\rm i} e_{i_1} \wedge \cdots \wedge e_{i_k}.
\end{aligned}
\]
Hence
\[
\omega_\C \cdot e_{i_1} \wedge \cdots \wedge e_{i_k} = {\rm i}^n {\rm i}^k (-{\rm i})^{n-k} e_{i_1} \wedge \cdots \wedge e_{i_k} = (-1)^{2n-k} e_{i_1} \wedge \cdots \wedge e_{i_k},
\]
which completes the proof.
\end{proof}

\begin{Remark}
\label{lift}
For $(V,J)$ an oriented $2n$-dimensional Euclidean real vector space with an orthogonal complex structure $J$ compatible with the orientation, the line in $S^+_\C$ corresponding to $J\in J_+(V)$ is $\Lambda^0_\C(V)\subset S_\C^+$.
Indeed,
\[(v \otimes 1 + w \otimes {\rm i}) \cdot 1 = v+Jw = 0 \ \Leftrightarrow \ w=Jv. \]
\end{Remark}

